\section{A Verified Compiler for Adaptive Oracle Programs}
\label{sec:trace-compiler}

The rules above compose a sequence of sampling operations.
An oracle program does not provide such a sequence: the location and arguments of its next call may depend on previous answers.
Our compiler is the verified adapter between these views.
It exposes a bounded number of calls to one selected operation without
restricting the surrounding program's control flow, and it comes with both a
same-package correctness theorem and a generic quantitative replacement theorem.

\subsection{A trace-based selected-call transformation}

Given a program \(A\) and a chosen operation \(p\), a paper proof can simply
say: run \(A\) until its next call to \(p\), answer that call, and resume.
In SSProve, the code after a call is represented by a
Rocq function from the answer to the rest of the program.  Such arbitrary
functions cannot be serialized into SSProve program data.  The compiler therefore records a
trace of the intervening calls, heap operations, and samples.
When it reaches \(p\), it returns the query without executing the call.  The
driver obtains an answer from a chosen implementation \(P_s\), adds it to the
trace, replays the recorded results through \(A\) to recover the code after the
call, and repeats this process up to \(q\) times.

Write \(\mathsf{Expose}_q(A,p;P_s)\) for the raw-code transformation that
exposes the first \(q\) calls made by \(A\) to operation \(p\) and resolves
them with \(P_s\).  Its factoring phase is trace-valued, but replay restores
the residual program, so the complete transformation has the same result type
as \(A\).  Write \(\mathsf{Compile}_q(A,p;P_s,P_r)\) for
\(\mathsf{Expose}_q(A,p;P_s)\) linked with residual package \(P_r\).

\subsection{Same-package correctness}

The transformation inserts trace bookkeeping, so before linking its syntax
differs from that of \(A\).  If both the exposed calls and the residual calls
are implemented by \(P\), however, each call receives the same answer as in
the original linked program, and trace replay reconstructs the same remaining
code.  The compiler therefore preserves behavior in the same-package case.

\begin{theorem}[Same-package compilation]
\label{thm:compile-same}
Subject to code and package validity and the presence of the selected
operation, for every initial heap,
\[
 \mathsf{Compile}_q(A,p;P,P)
 \quad\text{and}\quad
 \mathsf{Link}(A,P)
\]
have identical SSProve output/heap semantics.
\end{theorem}

The proof is an induction over the selected-call traversal.  At each step, the
recorded result is supplied to the corresponding continuation of the original
program, while validity rules out malformed traces.  Consequently, the
compiled and original programs induce the same output--heap distribution.

\subsection{A generic local-to-adaptive replacement theorem}

The following theorem is the main program-logic bridge.

\begin{theorem}[Replace $q$ calls]
\label{thm:compile-replace}
Let \(P_I\) synchronize equal operation inputs and equal heaps satisfying
\(I\).  Suppose two implementations \(P,P'\) of an operation \(p\) satisfy
the one-call Pythagorean judgment
\[
 \vDash\Pyth\{P_I\}\ P.p\approx_s P'.p\ \{I\}
\]
for a nonnegative tuple \(s\), preserve \(I\), and meet the compiler's typing,
footprint, and memory-separation conditions.
For every concrete SSProve program \(A\), replacing its first
\(q\) selected calls yields
\[
\begin{aligned}
 &\vDash\AEJ\{P_I\}\ \mathsf{Compile}_q(A,p;P,P)\\
 &\quad\approx_{\sqrt{q\norm{s}_1/2}}
   \mathsf{Compile}_q(A,p;P',P)\ \{o_L=o_R\},
\end{aligned}
\]
\end{theorem}
\begin{proof}[Proof sketch]
We first show that \(\mathsf{Compile}_q(A,p;P,P)\) and
\(\mathsf{Compile}_q(A,p;P',P)\) satisfy the expected Pythagorean judgment with
\(q\) concatenated copies of the tuple \(s\), by induction on \(q\) and the
Pythagorean sequence rule.
Each exposed call contributes \(s\) and preserves the invariant \(I\).

After \(q\) calls, the accumulated tuple has norm
\(q\norm{s}_1\); \textsc{Micciancio--Walter} is applied once, giving the
desired additive error bound.
\end{proof}

The theorem is generic in the program, packages, selected oracle, invariant, and budget tuple.
The noise-flooding game supplies only the one-call premise proved in
\Cref{sec:closing-hop}; the program logic and compiler then handle adversarial
control flow.
Combined with \Cref{thm:compile-same}, this is our quantitative counterpart of SSProve's local relational reasoning for package indistinguishability.

\subsection{Instantiating the compiler for noise flooding}
\label{sec:closing-hop}

We now instantiate the generic machinery to prove the missing arrow in
\Cref{eq:conceptual-game-chain}.
Let \(C_{\cA}^{\mathsf{open}}\) be the adversary-facing IND-CPAD game code
before its oracle implementations are linked; its encryption, evaluation, and
decryption calls remain unresolved.
Let \(P_{\mathsf{real}}\) implement the real encryption, evaluation, and
flooded-decryption oracles, and let \(P_{\mathsf{sim}}\) differ only by
flooding around the common recorded plaintext when the row's two plaintexts
agree.

The selected decryption body factors into a deterministic prefix, one Gaussian
sample, and deterministic postprocessing.
The invariant \(\Phi\) from \Cref{sec:adaptive-invariant} gives
\(d(\mathsf{dec}_0(sk,c),m)\leq e\) on every reachable row whose recorded
plaintexts agree.
Translation of charts by \Cref{eq:charts} puts the two answer distributions
over the same coordinate space.
Applying \Cref{eq:dg-kl} in \(n\) coordinates with
\(\sigma_e=\max\{1,e\}\gamma\) yields the following local judgment.

\begin{lemma}[One-call replacement]
\label{lem:one-call}
For any decryption input and paired heaps satisfying \(\Phi\), the real and
simulated decryption bodies satisfy a Pythagorean relational judgment with
budget tuple
\[
 (0,\epsilon_{\mathsf{nf}},0),
 \qquad
 \epsilon_{\mathsf{nf}}=\frac{n}{2\gamma^2},
\]
and restore \(\Phi\).
\end{lemma}

The zero entries record the shared prefix and postprocessing.
Branches with differing recorded plaintexts, as well as assertion-failing
branches, agree exactly.
Applying \Cref{thm:compile-replace} to \Cref{lem:one-call} concatenates one
such tuple for each of the first \(q\) decryption calls, preserving all
dependencies between queries.
The \textsc{Micciancio--Walter} rule is then invoked once, after the accumulated
KL budget is \(q\epsilon_{\mathsf{nf}}\), proving
\[
 \TVD(G_0,G_1)\leq\sqrt{q\epsilon_{\mathsf{nf}}/2}.
\]
Data processing projects completed output/heap pairs to the returned Boolean.
This proves \Cref{eq:central-game-hop}.

\smallskip\noindent\emph{From games to program representations.}
The \(G_i\) in \Cref{eq:conceptual-game-chain} are complete experiments: the
adversary has already been linked to its oracle implementations, so no
unresolved calls remain.
This is the usual cryptographic view of a game, but it hides the program
structure needed by \Cref{thm:compile-replace}.
The formal proof therefore makes the linking and call-compilation steps
explicit.
For selected-call implementation \(P_s\) and residual implementation \(P_r\),
write
\[
 H(P_s,P_r)=
 \mathsf{Compile}_{q}
   (C_{\cA}^{\mathsf{open}},\mathsf{ODec};P_s,P_r).
\]
Then
\begin{align*}
 H_0 &= G_0,\\
 H_1 &= \mathsf{Link}(C_{\cA}^{\mathsf{open}},P_{\mathsf{real}}),\\
 H_2 &= H(P_{\mathsf{real}},P_{\mathsf{real}}),\\
 H_3 &= H(P_{\mathsf{sim}},P_{\mathsf{real}}),\\
 H_4 &= H(P_{\mathsf{sim}},P_{\mathsf{sim}}),\\
 H_5 &= \mathsf{Link}(C_{\cA}^{\mathsf{open}},P_{\mathsf{sim}}),\\
 H_6 &= G_2.
\end{align*}
This is an implementation-level factorization of the two arrows in
\Cref{eq:conceptual-game-chain}, not seven cryptographic hybrids:
\(H_0,H_1,H_2\) are real-oracle representations of \(G_0\);
\(H_3,H_4,H_5\) are simulated-oracle representations of \(G_1\); and
\(H_6\) is the reduction game \(G_2\).
Rocq verifies
\begin{equation}
\label{eq:hybrid-chain}
 H_0\xrightarrow{0}H_1\xrightarrow{0}H_2
 \xrightarrow{\sqrt{q\epsilon_{\mathsf{nf}}/2}}H_3
 \xrightarrow{0}H_4\xrightarrow{0}H_5\xrightarrow{0}H_6.
\end{equation}
Here \(H_0=H_1\) unfolds package linking, and \(H_1=H_2\) is same-package
compiler correctness.
The single paid step \(H_2\to H_3\) is the application above.
For \(H_3=H_4\), the shared counter ensures that any residual decryption call
after the \(q\) selected entries assertion-fails on both sides.
The equality \(H_4=H_5\) removes the compiler using
\Cref{thm:compile-same}.
Finally, \(H_5\to H_6\) relates the simulated IND-CPAD presentation to the
concrete IND-CPA reduction by an exact heap invariant; the artifact factors
this relation through auxiliary linked programs.

Collapsing the zero-cost equalities in \Cref{eq:hybrid-chain} recovers the
three-game chain in \Cref{eq:conceptual-game-chain}; only \(H_2\to H_3\)
changes a distribution.
Triangle composition of \Cref{eq:hybrid-chain}, followed by the assumed
IND-CPA bound at \(H_6\), proves \Cref{thm:main}.
